\chapter{Angle inscrit, angle au centre}

\noindent\textit{Dans un cercle, deux types d'angles permettent de « voir » un même arc :
l'angle au centre, dont le sommet est le centre du cercle, et l'angle inscrit, dont le
sommet est un point du cercle. Une relation simple les relie, avec une conséquence
remarquable : tout triangle inscrit dans un cercle avec un côté sur un diamètre est
rectangle.}
\vspace{8pt}

\connecteBanner

\section{Angle inscrit, angle au centre}

\begin{definition}[Angle au centre, angle inscrit]
Soit un cercle de centre $O$ passant par $A$ et $B$. L'\textbf{angle au centre} $\widehat{AOB}$
a pour sommet le centre $O$. Pour un point $C$ du cercle (différent de $A$ et $B$),
l'\textbf{angle inscrit} $\widehat{ACB}$ a pour sommet $C$, sur le cercle. Les deux angles
\textbf{interceptent} le même arc $AB$ (celui qui ne contient pas $C$).
\end{definition}

\begin{illus}
\begin{tikzpicture}[scale=0.9]
\draw[onbline,line width=0.8pt] (0,0) circle (3);
\coordinate (O) at (0,0);
\coordinate (A) at (-1.93,2.30);
\coordinate (B) at (1.93,2.30);
\coordinate (C) at (0,-3);
\draw[onbo,line width=1.2pt] (O) -- (A) (O) -- (B);
\draw[onbblue,line width=1.2pt] (C) -- (A) (C) -- (B);
\fill (O) circle(1.6pt) node[below] {$O$};
\fill (A) circle(1.6pt) node[above] {$A$};
\fill (B) circle(1.6pt) node[above] {$B$};
\fill (C) circle(1.6pt) node[below] {$C$};
\pic[draw=onbo,angle radius=0.6cm,"\small $80^{\circ}$"] {angle=B--O--A};
\pic[draw=onbblue,angle radius=0.7cm,"\small $40^{\circ}$"] {angle=B--C--A};
\end{tikzpicture}
\fig{Figure — Angle au centre $\widehat{AOB}=80^{\circ}$, angle inscrit
$\widehat{ACB}=40^{\circ}$.}
\end{illus}

\section{Propriété : angle au centre et angle inscrit}

\begin{propriete}[Angle au centre = 2 $\times$ angle inscrit]
Si $\widehat{AOB}$ est l'angle au centre et $\widehat{ACB}$ l'angle inscrit interceptant le
même arc $AB$, alors
\[
\widehat{AOB}=2\times\widehat{ACB}.
\]
\end{propriete}

\begin{illus}
\begin{tikzpicture}[every node/.style={font=\small}]
\node[draw=onbblue,fill=white,rounded corners=3pt,minimum width=4.4cm,minimum height=0.7cm] (c1) {angle inscrit $\widehat{ACB}$};
\node[draw=onbo,fill=onbsoft,rounded corners=3pt,minimum width=4.4cm,minimum height=0.7cm,below=8mm of c1] (c2) {angle au centre $\widehat{AOB}=2\times\widehat{ACB}$};
\draw[->,onbmuted,line width=1pt] (c1) -- (c2);
\end{tikzpicture}
\fig{Figure — L'angle au centre vaut le double de l'angle inscrit qui intercepte le même
arc.}
\end{illus}

\begin{exemple}
Si $\widehat{ACB}=35^{\circ}$, alors $\widehat{AOB}=2\times35^{\circ}=70^{\circ}$.
Inversement, si $\widehat{AOB}=94^{\circ}$, alors $\widehat{ACB}=\dfrac{94^{\circ}}2=47^{\circ}$.
\end{exemple}

\section{Angles inscrits interceptant le même arc}

\begin{propriete}[Égalité des angles inscrits]
Deux angles inscrits qui interceptent le \textbf{même arc} (sommets du même côté de cet
arc) sont \textbf{égaux}.
\end{propriete}

\begin{illus}
\begin{tikzpicture}[scale=0.9]
\draw[onbline,line width=0.8pt] (0,0) circle (3);
\coordinate (A) at (-1.93,2.30);
\coordinate (B) at (1.93,2.30);
\coordinate (C1) at (-1.03,-2.82);
\coordinate (C2) at (2.30,-1.93);
\draw[onbo,line width=1.1pt] (C1) -- (A) (C1) -- (B);
\draw[onbblue,line width=1.1pt] (C2) -- (A) (C2) -- (B);
\fill (A) circle(1.6pt) node[above] {$A$};
\fill (B) circle(1.6pt) node[above] {$B$};
\fill (C1) circle(1.6pt) node[below left] {$C_1$};
\fill (C2) circle(1.6pt) node[below right] {$C_2$};
\pic[draw=onbo,angle radius=0.6cm,"\small $40^{\circ}$"] {angle=B--C1--A};
\pic[draw=onbblue,angle radius=0.6cm,"\small $40^{\circ}$"] {angle=B--C2--A};
\end{tikzpicture}
\fig{Figure — $\widehat{AC_1B}$ et $\widehat{AC_2B}$ interceptent le même arc : ils sont
égaux.}
\end{illus}

\begin{exemple}
$\widehat{AC_1B}=47^{\circ}$ et $C_2$ intercepte le même arc $AB$ que $C_1$ : alors
$\widehat{AC_2B}=47^{\circ}$ également.
\end{exemple}

\section{Triangle rectangle inscrit dans un cercle}

\begin{propriete}[Cas du diamètre]
Si $[AB]$ est un \textbf{diamètre} d'un cercle et $C$ un point de ce cercle (distinct de
$A$ et $B$), alors le triangle $ABC$ est \textbf{rectangle en $C$}.
\end{propriete}

\begin{illus}
\begin{tikzpicture}[scale=0.9]
\draw[onbline,line width=0.8pt] (0,0) circle (3);
\coordinate (A) at (-3,0);
\coordinate (B) at (3,0);
\coordinate (C) at (1.5,2.598);
\draw[onbline,line width=0.8pt] (A) -- (B);
\draw[onbo,line width=1.2pt] (A) -- (C) -- (B);
\fill (A) circle(1.6pt) node[left] {$A$};
\fill (B) circle(1.6pt) node[right] {$B$};
\fill (C) circle(1.6pt) node[above] {$C$};
\fill (0,0) circle(1.6pt) node[below] {$O$};
\pic[draw=onbblue,angle radius=0.5cm,"\small $90^{\circ}$"] {angle=A--C--B};
\end{tikzpicture}
\fig{Figure — $AB$ diamètre : le triangle $ABC$ est rectangle en $C$, quel que soit $C$ sur
le cercle.}
\end{illus}

\begin{remarque}[Justification]
$\widehat{AOB}=180^{\circ}$ (angle plat, $O$ est le milieu de $[AB]$), donc l'angle inscrit
$\widehat{ACB}=\dfrac{180^{\circ}}2=90^{\circ}$ : c'est un cas particulier de la propriété
précédente.
\end{remarque}

\begin{propriete}[Réciproque]
Si un triangle $ABC$ est rectangle en $C$, alors $C$ appartient au cercle de diamètre
$[AB]$ (le centre de ce cercle est le milieu de $[AB]$).
\end{propriete}

\begin{exemple}
$AB$ est un diamètre de $15$ cm, $C$ un point du cercle avec $AC=9$ cm. Le triangle $ABC$
est rectangle en $C$, donc (Pythagore) $BC^{2}=AB^{2}-AC^{2}=225-81=144$, soit $BC=12$ cm.
\end{exemple}

% ═══════════════════════════════════════════════════════════════
\section{L'essentiel à retenir}

\begin{synthese}
\textbf{Angle au centre / inscrit.} Même arc intercepté : $\widehat{AOB}=2\times\widehat{ACB}$.\\[4pt]
\textbf{Même arc, plusieurs sommets.} Les angles inscrits interceptant le même arc sont
égaux.\\[4pt]
\textbf{Diamètre.} $[AB]$ diamètre, $C$ sur le cercle $\Rightarrow ABC$ rectangle en $C$
(et réciproquement).\\[4pt]
\textbf{Triangle $OAB$.} Isocèle en $O$ (car $OA=OB=$ rayon) : ses deux angles à la base
sont égaux.\\[4pt]
\textbf{Piège fréquent.} Vérifier que les angles interceptent bien le \textbf{même} arc
avant d'appliquer l'égalité ; ne pas confondre angle au centre et angle inscrit dans la
relation « double ».
\end{synthese}

% ═══════════════════════════════════════════════════════════════
\section{Méthodes \& savoir-faire}

\methode{Calculer un angle au centre connaissant l'angle inscrit}
Multiplier l'angle inscrit par $2$.
\begin{exemple}
$\widehat{ACB}=52^{\circ}$ : $\widehat{AOB}=2\times52^{\circ}=104^{\circ}$.
\end{exemple}

\methode{Calculer un angle inscrit connaissant l'angle au centre}
Diviser l'angle au centre par $2$.
\begin{exemple}
$\widehat{AOB}=136^{\circ}$ : $\widehat{ACB}=\dfrac{136^{\circ}}2=68^{\circ}$.
\end{exemple}

\methode{Utiliser l'égalité des angles inscrits interceptant le même arc}
Vérifier que les deux angles inscrits ont pour côtés les mêmes points $A$, $B$ et
interceptent le même arc, puis conclure qu'ils sont égaux.
\begin{exemple}
$\widehat{ACB}=63^{\circ}$, $D$ sur le même arc que $C$ : $\widehat{ADB}=63^{\circ}$.
\end{exemple}

\methode{Reconnaître un triangle rectangle inscrit dans un cercle}
Vérifier qu'un des côtés du triangle est un \textbf{diamètre} du cercle circonscrit : le
sommet opposé donne alors l'angle droit.
\begin{exemple}
$AB$ diamètre, $C$ sur le cercle : le triangle $ABC$ est rectangle en $C$, sans calcul
supplémentaire.
\end{exemple}

\methode{Utiliser la réciproque}
Si un triangle est rectangle, son cercle circonscrit a pour diamètre l'hypoténuse (le
centre est le milieu de l'hypoténuse).
\begin{exemple}
Triangle $ABC$ rectangle en $C$, $AB=10$ cm : le rayon du cercle circonscrit est
$\dfrac{10}2=5$ cm.
\end{exemple}

\methode{Résoudre un problème combinant plusieurs propriétés}
Identifier d'abord le diamètre (angle droit), puis utiliser Pythagore ou les angles au
centre/inscrits selon la question posée.
\begin{exemple}
$BC$ diamètre $=26$ cm, $AB=10$ cm : le triangle est rectangle en $A$, donc
$AC=\sqrt{26^{2}-10^{2}}=\sqrt{576}=24$ cm.
\end{exemple}

\methode{Calculer les angles du triangle isocèle $OAB$}
Utiliser $OA=OB=$ rayon (triangle isocèle en $O$) : les deux angles à la base valent
$\dfrac{180^{\circ}-\widehat{AOB}}2$.
\begin{exemple}
$\widehat{AOB}=100^{\circ}$ : les angles en $A$ et $B$ du triangle $OAB$ valent chacun
$\dfrac{180^{\circ}-100^{\circ}}2=40^{\circ}$.
\end{exemple}

% ═══════════════════════════════════════════════════════════════
\section{Exercices d'application}
\begin{multicols}{2}

\rubrique{Angle au centre et angle inscrit}
\exo{}\diff{1} L'angle inscrit $\widehat{ACB}=35^{\circ}$. Calculer l'angle au centre
$\widehat{AOB}$.

\exo{}\diff{1} L'angle au centre $\widehat{AOB}=94^{\circ}$. Calculer l'angle inscrit
$\widehat{ACB}$.

\exo{}\diff{2} L'angle inscrit $\widehat{ACB}=52^{\circ}$. Calculer l'angle au centre
$\widehat{AOB}$, puis un autre angle inscrit $\widehat{ADB}$ interceptant le même arc.

\exo{}\diff{2} L'angle au centre $\widehat{AOB}=136^{\circ}$. Calculer l'angle inscrit
$\widehat{ACB}$.

\rubrique{Angles inscrits interceptant le même arc}
\exo{}\diff{1} $\widehat{ACB}=47^{\circ}$, $D$ intercepte le même arc $AB$ que $C$.
Calculer $\widehat{ADB}$.

\exo{}\diff{2} $C$, $D$, $E$ sont trois points du même arc, avec $\widehat{ACB}=63^{\circ}$.
Calculer $\widehat{ADB}$ et $\widehat{AEB}$.

\exo{}\diff{2} $\widehat{ACB}=38^{\circ}$. Calculer l'angle au centre $\widehat{AOB}$, puis
un autre angle inscrit $\widehat{ADB}$ interceptant le même arc.

\exo{}\diff{2} Sachant que l'angle au centre vaut le double de l'angle inscrit, et que deux
angles inscrits interceptant le même arc valent chacun $29^{\circ}$, calculer l'angle au
centre correspondant.

\rubrique{Triangle rectangle inscrit (diamètre)}
\exo{}\diff{1} $AB$ est un diamètre d'un cercle, $C$ un point du cercle. Que peut-on dire
de l'angle $\widehat{ACB}$ ? Justifier.

\exo{}\diff{2} Un triangle $ABC$ rectangle en $C$ a pour hypoténuse $AB=10$ cm. Quel est le
rayon de son cercle circonscrit ?

\exo{}\diff{2} $AB$ est un diamètre de $13$ cm, $C$ un point du cercle avec $AC=5$ cm.
Calculer $BC$.

\exo{}\diff{2} $BC$ est un diamètre de $17$ cm, $A$ un point du cercle avec $AB=8$ cm.
Calculer $AC$.

\rubrique{Problèmes combinés}
\exo{}\diff{2} Un triangle $ABC$ rectangle en $A$ est inscrit dans un cercle de diamètre
$BC=26$ cm, avec $AB=10$ cm. Calculer $AC$ et le rayon du cercle.

\exo{}\diff{2} Dans un cercle de centre $O$, l'angle au centre $\widehat{AOB}=110^{\circ}$.
Calculer l'angle inscrit $\widehat{ACB}$, puis les deux autres angles du triangle isocèle
$OAB$.

\exo{}\diff{2} Une fenêtre semi-circulaire a pour diamètre $AB=2$ m. Un point $C$ sur l'arc
forme un triangle $ABC$. Quelle est la mesure de l'angle $\widehat{ACB}$ ?

\exo{}\diff{2} Dans un cercle de centre $O$, l'angle inscrit $\widehat{ACB}=41^{\circ}$ et
$D$ intercepte le même arc $AB$ que $C$. Calculer $\widehat{ADB}$ et l'angle au centre
$\widehat{AOB}$.

% ═══════════════════════════════════════════════════════════════
\end{multicols}

\section{Exercices d'entraînement}
\begin{multicols}{2}

\rubrique{Angle au centre et angle inscrit}
\exo{}\diff{1} $\widehat{ACB}=28^{\circ}$. Calculer $\widehat{AOB}$.

\exo{}\diff{1} $\widehat{AOB}=150^{\circ}$. Calculer $\widehat{ACB}$.

\exo{}\diff{2} $\widehat{ACB}=61^{\circ}$. Calculer $\widehat{AOB}$.

\exo{}\diff{2} $\widehat{AOB}=88^{\circ}$. Calculer $\widehat{ACB}$.

\rubrique{Angles interceptant le même arc}
\exo{}\diff{2} $\widehat{ACB}=55^{\circ}$, $D$ intercepte le même arc. Calculer
$\widehat{ADB}$.

\exo{}\diff{2} $\widehat{ACB}=72^{\circ}$, $E$ intercepte le même arc. Calculer
$\widehat{AEB}$.

\exo{}\diff{2} L'angle au centre $\widehat{AOB}=100^{\circ}$. Calculer les deux angles
inscrits égaux $\widehat{ACB}$ et $\widehat{ADB}$ interceptant le même arc.

\exo{}\diff{2} Deux angles inscrits interceptant le même arc valent chacun $39^{\circ}$.
Justifier qu'ils sont égaux et calculer l'angle au centre correspondant.

\rubrique{Triangle rectangle inscrit}
\exo{}\diff{2} $AB$ diamètre $=15$ cm, $C$ sur le cercle avec $AC=9$ cm. Calculer $BC$.

\exo{}\diff{2} $AB$ diamètre $=25$ cm, $C$ sur le cercle avec $AC=7$ cm. Calculer $BC$.

\exo{}\diff{2} $AB$ diamètre $=20$ cm, $C$ sur le cercle avec $AC=12$ cm. Calculer $BC$.

\exo{}\diff{2} Un triangle rectangle a pour hypoténuse $13$ cm. Quel est le rayon de son
cercle circonscrit ? Si un côté de l'angle droit mesure $5$ cm, calculer l'autre.

\rubrique{Triangle isocèle $OAB$}
\exo{}\diff{2} $\widehat{AOB}=120^{\circ}$. Calculer les deux autres angles du triangle
isocèle $OAB$.

\exo{}\diff{2} $\widehat{AOB}=70^{\circ}$. Calculer les deux autres angles du triangle
isocèle $OAB$.

\exo{}\diff{2} $\widehat{AOB}=88^{\circ}$. Calculer les deux autres angles du triangle
isocèle $OAB$.

\exo{}\diff{2} $\widehat{AOB}=46^{\circ}$. Calculer les deux autres angles du triangle
isocèle $OAB$.

\rubrique{Problèmes}
\exo{}\diff{2} Un triangle $ABC$ rectangle en $B$ est inscrit dans un cercle de diamètre
$AC=41$ cm, avec $AB=9$ cm. Calculer $BC$.

\exo{}\diff{2} Une place circulaire de rayon $6{,}5$ m a un diamètre $AB$. Un lampadaire
est installé en un point $C$ du cercle avec $AC=5$ m. Calculer $BC$.

\exo{}\diff{2} Dans un cercle de centre $O$, l'angle inscrit $\widehat{ACB}=44^{\circ}$.
Calculer l'angle au centre $\widehat{AOB}$, puis les deux angles à la base du triangle
isocèle $OAB$.

\exo{}\diff{3} Un pont en arc de cercle a pour diamètre (à la base) $AB=30$ m. Le sommet
$C$ de l'arc vérifie $AC=18$ m. Calculer $BC$.

% ═══════════════════════════════════════════════════════════════
\end{multicols}

\section{Sujets type BEPC}

\sujetbac[angle au centre/inscrit et triangle rectangle inscrit]
\begin{enumerate}[label=\arabic*)]
\item Dans un cercle de centre $O$, l'angle inscrit $\widehat{ACB}=42^{\circ}$. Calculer
l'angle au centre $\widehat{AOB}$.
\item $D$ est un autre point du cercle tel que $\widehat{ADB}$ intercepte le même arc $AB$
que $\widehat{ACB}$. Calculer $\widehat{ADB}$.
\item $AB$ est un diamètre du cercle, $C$ un point du cercle avec $AC=9$ cm et $AB=15$ cm.
Justifier que le triangle $ABC$ est rectangle en $C$, puis calculer $BC$.
\end{enumerate}

\sujetbac[triangle isocèle et cercle circonscrit]
\begin{enumerate}[label=\arabic*)]
\item Dans un cercle de centre $O$ et de rayon $6{,}5$ cm, l'angle au centre
$\widehat{AOB}=100^{\circ}$. Calculer l'angle inscrit $\widehat{ACB}$ ($C$ sur le grand
arc).
\item Le triangle $OAB$ est isocèle en $O$. Calculer ses deux autres angles.
\item $BC$ est un diamètre du cercle, mesurant $13$ cm. Un point $A$ du cercle vérifie
$AB=5$ cm. Calculer $AC$.
\end{enumerate}

\sujetbac[angles et fenêtre semi-circulaire]
\begin{enumerate}[label=\arabic*)]
\item Dans un cercle, deux angles inscrits $\widehat{ACB}$ et $\widehat{ADB}$ interceptent
le même arc $AB$, avec $\widehat{ACB}=58^{\circ}$. Calculer $\widehat{ADB}$ et l'angle au
centre $\widehat{AOB}$.
\item Une fenêtre semi-circulaire a pour diamètre $AB=2{,}6$ m. Un point $C$ sur l'arc
forme un triangle $ABC$ avec $AC=1$ m. Justifier que le triangle est rectangle en $C$, puis
calculer $BC$.
\item Le triangle $OAB$ est isocèle en $O$, avec $\widehat{AOB}=116^{\circ}$. Calculer les
deux autres angles.
\end{enumerate}

% ═══════════════════════════════════════════════════════════════
\section{Corrigés}
\begin{multicols}{2}

\corrige{de l'exercice 1}
$\widehat{AOB}=2\times35^{\circ}=70^{\circ}$.

\corrige{de l'exercice 2}
$\widehat{ACB}=\dfrac{94^{\circ}}2=47^{\circ}$.

\corrige{de l'exercice 3}
$\widehat{AOB}=2\times52^{\circ}=104^{\circ}$ ; $\widehat{ADB}=52^{\circ}$ (même arc).

\corrige{de l'exercice 4}
$\widehat{ACB}=\dfrac{136^{\circ}}2=68^{\circ}$.

\corrige{de l'exercice 5}
$\widehat{ADB}=47^{\circ}$ (même arc intercepté).

\corrige{de l'exercice 6}
$\widehat{ADB}=\widehat{AEB}=63^{\circ}$.

\corrige{de l'exercice 7}
$\widehat{AOB}=2\times38^{\circ}=76^{\circ}$ ; $\widehat{ADB}=38^{\circ}$.

\corrige{de l'exercice 8}
$\widehat{AOB}=2\times29^{\circ}=58^{\circ}$.

\corrige{de l'exercice 9}
$\widehat{ACB}=90^{\circ}$ : un côté du triangle est un diamètre, donc le triangle est
rectangle au sommet opposé.

\corrige{de l'exercice 10}
Rayon $=\dfrac{10}2=5$ cm.

\corrige{de l'exercice 11}
Triangle rectangle en $C$ : $BC^{2}=13^{2}-5^{2}=144$, donc $BC=12$ cm.

\corrige{de l'exercice 12}
Triangle rectangle en $A$ : $AC^{2}=17^{2}-8^{2}=225$, donc $AC=15$ cm.

\corrige{de l'exercice 13}
Triangle rectangle en $A$ : $AC^{2}=26^{2}-10^{2}=576$, donc $AC=24$ cm. Rayon
$=\dfrac{26}2=13$ cm.

\corrige{de l'exercice 14}
$\widehat{ACB}=\dfrac{110^{\circ}}2=55^{\circ}$. Triangle $OAB$ isocèle en $O$ : angles à
la base $=\dfrac{180^{\circ}-110^{\circ}}2=35^{\circ}$ chacun.

\corrige{de l'exercice 15}
$\widehat{ACB}=90^{\circ}$ (angle inscrit interceptant un diamètre).

\corrige{de l'exercice 16}
$\widehat{ADB}=41^{\circ}$ (même arc) ; $\widehat{AOB}=2\times41^{\circ}=82^{\circ}$.

\corrige{du sujet 1}
1) $\widehat{AOB}=2\times42^{\circ}=84^{\circ}$.\\
2) $\widehat{ADB}=42^{\circ}$ (même arc).\\
3) $[AB]$ diamètre $\Rightarrow$ rectangle en $C$ ; $BC^{2}=15^{2}-9^{2}=144$, donc
$BC=12$ cm.

\corrige{du sujet 2}
1) $\widehat{ACB}=\dfrac{100^{\circ}}2=50^{\circ}$.\\
2) Angles à la base $=\dfrac{180^{\circ}-100^{\circ}}2=40^{\circ}$ chacun.\\
3) Rectangle en $A$ ($BC$ diamètre) : $AC^{2}=13^{2}-5^{2}=144$, donc $AC=12$ cm.

\corrige{du sujet 3}
1) $\widehat{ADB}=58^{\circ}$ ; $\widehat{AOB}=2\times58^{\circ}=116^{\circ}$.\\
2) $[AB]$ diamètre $\Rightarrow$ rectangle en $C$ ;
$BC^{2}=2{,}6^{2}-1^{2}=5{,}76$, donc $BC=2{,}4$ m.\\
3) Angles à la base $=\dfrac{180^{\circ}-116^{\circ}}2=32^{\circ}$ chacun.

\end{multicols}
\exercicesBanner
